\begin{textsample}{POS Dim 1 – am – Score 38.00 – am/am\_ef\_19.txt}  \label{ex:f1_pos_216}
6.
A formação do professor A Educação Especial na perspectiva da Educação Inclusiva é \textbf{uma} temática bem discutida na sociedade \textbf{contemporânea}.
Reforça e objetiva incluir todos no âmbito escolar e na sociedade e, com isso, fazer cumprir o \textbf{que} garante a Constituição Federal Brasileira, \textbf{que} afirma \textbf{que} todos têm direito à educação.
Torna -se indispensável \textbf{que} a formação do professor da Educação Especial na perspectiva da Educação Inclusiva seja \textbf{uma} \textbf{abordagem} objetiva e reflita sobre \textbf{uma} melhor compreensão dessa modalidade de ensino. \textbf{Uma} forma de compreender melhor esse modelo de ensino é conhecer \textbf{um} pouco da história e das legislações \textbf{que} tratam sobre essa modalidade.
Identificar também como é realizada a formação do professor \textbf{que} atua de forma direta nesse viés, buscando encontrar em sua formação teórico-prática elementos \textbf{que} proporcionem a inclusão do público-alvo da Educação Especial nas escolas de ensino regular e, consequentemente, na sociedade. 6.1. \textbf{Historicamente}, as pessoas com deficiência sofreram com várias formas de preconceito \textbf{que} se estenderam durante muitos \textbf{séculos} e, apesar de \textbf{hoje} existirem leis \textbf{que} \textbf{amparam} seus direitos \textbf{enquanto} cidadãos, ainda se faz necessário \textbf{que} outras medidas sejam tomadas para impedir ou até amenizar algumas formas de \textbf{exclusão} \textbf{que} ainda são \textbf{imperativas} em nossa sociedade.
Com os avanços dos Direitos Humanos, registraram -se consideráveis progressos na conquista da igualdade e do exercício de direitos.
É o \textbf{que} se sente e o \textbf{que} se observa atualmente, tendo como grande enfoque a busca da inclusão destas pessoas \textbf{historicamente} marcadas pela segregação, pelo preconceito e pela rejeição.
Nesse processo, é de \textbf{suma} importância \textbf{que} haja \textbf{um} trabalho conjunto com toda a comunidade escolar, fazendo com \textbf{que} estejam todos cientes e envolvidos nesta causa, tornando o ambiente acessível para essas pessoas. \textbf{Um} ambiente educacional acessível e favorável a esses educandos \textbf{configura} \textbf{um} indicativo de \textbf{que} a inclusão pode de fato acontecer, e não somente a integração.
De acordo com Sassaki ( 2006 ), a integração propõe a inserção parcial do \textbf{sujeito}, \textbf{enquanto} \textbf{que} a inclusão propõe a inserção total.
Para \textbf{que} isso aconteça, a escola precisa romper com a perspectiva homogeneizadora e adotar estratégias para assegurar os direitos de aprendizagem de todos.
Segundo esse mesmo \textbf{autor}, as escolas precisam romper com diversas barreiras em vários aspectos de suas estruturas tradicionais para favorecer a acessibilidade no atendimento desses alunos.
De acordo com Sassaki ( 2009 ), são seis as principais barreiras a serem consideradas com total atenção : > [... ] arquitetônica ( sem barreiras físicas ), comunicacional ( sem barreiras na comunicação entre pessoas ), \textbf{metodológica} ( sem barreiras nos \textbf{métodos} e técnicas de lazer, trabalho, educação etc. ), instrumental ( sem barreiras nos instrumentos, ferramentas, utensílios etc. ), programática ( sem barreiras embutidas em políticas, legislações, normas etc. ) e atitudinal ( sem preconceitos, estereótipos, estigmas e discriminações nos comportamentos da sociedade para a pessoa \textbf{que} tem deficiência ). ( SASSAKI, 2009, p. 1-2 ) Nesse cenário, é \textbf{preciso} o envolvimento de vários profissionais de outras áreas de conhecimento, \textbf{que} atuarão em conjunto com os profissionais da educação, como os profissionais da saúde, fonoaudiólogos, psicólogos, assistentes sociais e, principalmente, a participação da família, por esta estabelecer \textbf{um} contato direto com a pessoa com deficiência.
Sendo assim, precisa estar bem informada a respeito da condição de seu assegurado e apta aos cuidados de seu familiar com deficiência, consolidando, assim, \textbf{um} trabalho interdisciplinar de qualidade e integral para com os diferentes.
O professor \textbf{que} atua de forma direta realiza seu trabalho pedagógico no ambiente escolar, busca promover as potencialidades dos \textbf{sujeitos} envolvidos no processo de \textbf{ensino-aprendizagem} e torna -se \textbf{um} profissional relevante e insubstituível nessa modalidade de educação.
A esse respeito, corrobora Cury ( 2003 ) : > Os educadores, apesar de suas dificuldades, são insubstituíveis porque a gentileza, a solidariedade, a tolerância, a inclusão, os sentimentos altruístas, enfim, todas as áreas da sensibilidade não podem ser ensinadas por máquinas, e sim por seres \textbf{humanos}.
A educação, como \textbf{um} fenômeno de responsabilidade social, possui \textbf{um} papel importante para o desenvolvimento e crescimento de \textbf{uma} nação \textbf{que} não pode ser ignorado por nenhum cidadão, sobretudo pelo poder público, pois é a escola \textbf{que} prepara os indivíduos para sua atuação e para exercer sua plena cidadania, gozando de todos os bens e serviços \textbf{que} sejam oferecidos por essa sociedade da qual faz parte.
Para \textbf{que} os resultados educacionais das pessoas com deficiência sejam alcançados com sucesso e, ainda, recebam \textbf{uma} formação adequada, faz -se necessário \textbf{que} os educadores envolvidos nesse processo, como gestores, pedagogos e professores, recebam \textbf{uma} formação de qualidade \textbf{que} possa contribuir de forma positiva para esses resultados e \textbf{que} estes atendam aos objetivos da sociedade.
A qualidade da formação desses profissionais, tanto inicial quanto continuada, é de total importância para refletir na qualidade do ensino de seus \textbf{discentes}.
A UNESCO, em seu relatório de 1998, relacionou a qualidade do ensino com a formação \textbf{qualificada} desses profissionais ( UNESCO, 1998 apud Roseli Rocha de Carvalho, Educação Especial : do querer ao fazer, p. 28 ) : > A qualidade dos serviços educacionais para pessoas com deficiência \textbf{depende} da qualidade da formação.
Esta deverá ser parte integrante dos planos nacionais, onde se contemplam os requisitos dessa formação.
Houve \textbf{um} momento na história do nosso país em \textbf{que} a formação para os professores era feita de \textbf{uma} forma técnica, apenas pelo magistério, em \textbf{que} o aluno do ensino médio era direcionado para escolas técnicas para receber \textbf{uma} formação direcionada a \textbf{uma} atuação como professor nas séries iniciais.
Os estudos para essa modalidade eram pautados nas metodologias de ensino, domínio dos conteúdos e dinâmica da turma, ou seja, nas práticas realizadas nas salas de aula.
Com a expansão do ensino no território nacional, surgiram vários problemas relacionados à qualidade do ensino e da aprendizagem.
Na década de noventa, \textbf{exigia} -se do profissional da educação \textbf{uma} nova postura, fundamentada na \textbf{teoria} construtivista interacionista concebida por Jean Piaget ( 1896-1980 ).
Para esse \textbf{teórico}, a criança, como \textbf{um} ser biológico, é ativa e age espontaneamente no meio, ou seja, é pelo contato com o mundo \textbf{que} seus conhecimentos são construídos.
Destaca Furtado ( 2009, p. 139 ) : > O ser humano, dotado de estrutura biológica, herdada por \textbf{uma} forma de funcionamento intelectual, ou seja, \textbf{uma} maneira de interagir com o ambiente \textbf{que} o leva à construção de \textbf{um} conjunto permitirá a organização desses significados em estruturas cognitivas.
É importante \textbf{que} o professor possua conhecimentos dessa natureza, \textbf{que} contempla a \textbf{psicologia} do desenvolvimento humano, disciplina essa \textbf{que} compõe a grade curricular do curso de formação inicial de professores.
Nesse sentido, para poder conhecer melhor seus educandos, através desse conhecimento e de \textbf{um} olhar cuidadoso, fica mais fácil para o professor identificar algumas deficiências ou dificuldades de aprendizagem, como dislexia, dislalia, deficiência intelectual, transtornos do espectro autista e outros.
Foi com o intuito de reparar e melhorar o atendimento educacional nas escolas brasileiras \textbf{que} houve \textbf{uma} exigência por parte do poder público vigente, a qual resultou na aprovação, em 2009, de \textbf{um} projeto \textbf{que} torna obrigatório \textbf{que} todos os professores do ensino básico tenham diploma universitário e licenciatura.
A Lei de Diretrizes e Bases da Educação ( LDB ) determina \textbf{que} todos os professores devam possuir curso superior.
A partir de então, o curso de \textbf{Pedagogia} passou a formar esses profissionais para atuar na Educação Infantil e nos Anos Iniciais.
São nessas etapas de ensino \textbf{que} inicialmente foram direcionados os estágios supervisionados, pois, neste primeiro momento, não constava nessa formação estudo para a Educação Especial.
Todavia, com a assinatura da Declaração de Salamanca, em 1994, o Brasil passou a rever o atendimento de crianças com deficiências, o \textbf{que} levantou \textbf{uma} questão a respeito da formação de professores, pois, no entendimento de alguns educadores, não era obrigatório formar professores especializados nessa modalidade de ensino, pois a inclusão passaria a ser \textbf{tarefa} de todas.
Diante dos desafios da Educação Especial, surge a importância do aperfeiçoamento dos professores para o exercício de serviços educacionais especializados.
A Portaria nº 1.793, de dezembro de 1994, recomendou a implementação, nos currículos de formação docente, da \textbf{disciplina} “ Aspectos ético-político educacionais da normatização e integração da pessoa com necessidades especiais ”.
Indicou a manutenção e expansão dos estudos como cursos e especialização na área desse conhecimento, com o objetivo de melhorar a qualidade do atendimento para esta modalidade.
Também consta na Lei de Diretrizes e Bases da Educação Nacional ( LDBN 9.394/96 ) \textbf{que} os professores tenham especialização adequada em nível médio ou superior para o atendimento especializado, bem como professores do ensino regular capacitados para a integração desses educandos nas classes comuns.
Consta no Plano Nacional de Educação ( PNE ), com vigência de 2014 a 2024, \textbf{que} seus eixos temáticos contemplem em suas diretrizes, metas e estratégias de ações \textbf{um} melhor atendimento nas escolas de todo o país : I ) Papel do Estado na garantia do direito à educação de qualidade : organização e regulamentação da educação nacional ; II ) Qualidade da educação, gestão democrática e avaliação ; III ) \textbf{Democratização} do acesso, permanência e sucesso escolar ; IV ) Formação e valorização dos/das profissionais da educação ; V ) Financiamento da educação e controle social ; VI ) Justiça social, educação e trabalho ; inclusão, diversidade e igualdade ( BRASIL, 2015 ).
Todos esses elementos legais vêm corroborando no sentido de tornar o sistema educacional brasileiro mais significativo e eficaz e, com isso, contribuindo com a qualidade da formação dos professores, conforme asseveram as metas 15 e 16 do referido documento : > [... ] assegurado \textbf{que} todos os professores e as professoras da educação básica possuam formação específica de nível superior, obtida em curso de licenciatura na área de conhecimento em \textbf{que} atuam. > > [... ] formar em nível de pós-graduação, 50 \% ( cinquenta por cento ) dos professores da educação básica, até o último ano de vigência deste PNE, e garantir a todos ( as ) os profissionais da educação básica formação continuada em sua área de atuação, considerando as necessidades, demanda e \textbf{contextualizações} de ensino. ( BRASIL, 2015, p. 79-80 ).
Entretanto, a educação das pessoas com deficiência, \textbf{que} no início estabelece o desenvolvimento de novas ações pedagógicas na intenção de explorar e potencializar habilidades e competências, foi mal interpretada.
Por falta de adequação ou incapacidade de seus educandos, passou a segregar mais do \textbf{que} incluir na interação social ( PACHECO ; ALVES, 2017 ).
Na tradição, o curso de \textbf{Pedagogia} não era destinado a formar professores, mas sim a formar educadores, planejadores, gestores e \textbf{pesquisadores} como bacharelado.
Depois, foi \textbf{agregado} mais \textbf{um} ano de didática e práticas de ensino para formar professor, mas não alfabetizar, e sim para dar aula nas escolas normais.
O processo de formação de professores ainda carece de maiores ajustes para se tornar \textbf{um} curso \textbf{que} forme de fato profissionais \textbf{que} atuem nas salas de aula e saibam desenvolver \textbf{um} trabalho junto a \textbf{um} público diversificado.
E por mais \textbf{que} ocorra a inserção de \textbf{métodos} singulares na construção da formação do educador, como já abordado, a \textbf{percepção} do docente pode não condizer com sua visão de Educação Especial trazida por essas práticas.
Por isso, é premente ressaltar o papel \textbf{que} cabe à formação continuada, \textbf{que} ocupou \textbf{um} espaço considerável, onde profissionais da educação têm a oportunidade de expandir e \textbf{agregar} novos conhecimentos, proporcionando ao docente em atividade \textbf{uma} excelente oportunidade para as trocas de experiência com outros profissionais da área, tornando -os mais \textbf{qualificados} para assumirem \textbf{uma} melhor postura diante de seus alunos.
Dessa maneira, infere Baumel ( 2003 ) : > A formação de professores deve ser concebida como \textbf{um} contínuo, associado à compreensão do desenvolvimento profissional ; em outras palavras, formar e articular “ \textbf{uma} variedade de formatos de aprendizagem ”.
O comprometimento, \textbf{aqui}, é de interligar a formação inicial com a continuada, \textbf{que} não abarca o termo e o processo de capacitação.
O processo de formação inicial e continuada é \textbf{um} projeto diferenciado, em fases, ao longo de \textbf{uma} finalidade e \textbf{um} estado de desenvolvimento profissional ( BAUMEL, 2003, p. 30 ).
Tratando -se de Educação Especial na perspectiva da educação inclusiva, os desafios são constantes e requerem desse profissional controle e domínio das diversas situações do cotidiano escolar voltadas a essa modalidade.
O professor pode encontrar essa formação no próprio espaço da escola e nos núcleos de apoio ofertados pelas redes de ensino nas quais estiver inserido, por meio de palestras, oficinas, leituras e cursos.
O mesmo deve ter a \textbf{vontade} de ir em busca de novos conhecimentos, no sentido de avançar com \textbf{um} currículo diferenciado em relação à qualidade de suas práticas educativas com seus educandos.
A esse respeito, aduz Prieto ( 2006 ) : > A formação continuada do professor deve ser \textbf{um} \textbf{compromisso} dos sistemas de ensino comprometidos com a qualidade do ensino \textbf{que}, nessa perspectiva, devem assegurar \textbf{que} sejam aptos a elaborar e a implantar novas propostas e práticas de ensino para responder às características de seus alunos, incluindo aquelas evidenciadas pelos alunos com necessidades educacionais especiais ( PRIETO, 2006, p. 57 ).
É por meio dessa formação continuada \textbf{que} o professor pode ampliar seu olhar para contemplar as potencialidades de seus alunos e não enxergar somente suas limitações.
Diante dos desafios \textbf{que} se referem à Educação Especial e Inclusiva, é indispensável \textbf{que} o professor mantenha \textbf{uma} constante formação teórico-prática, a fim de atender melhor nesse ambiente tão diversificado \textbf{que} é o chão das escolas.

% matched lemmas: abordagem, agregar, amparar, aqui, autor, compromisso, configurar, contemporâneo, contextualização, democratização, depender, discente, disciplina, enquanto, ensino-aprendizagem, exclusão, exigir, historicamente, hoje, humanos, imperativo, metodológico, método, pedagogia, percepção, pesquisador, preciso, psicologia, qualificar, que, sujeito, sumo, século, tarefa, teoria, teórico, um, vontade
\end{textsample}
